% ---------------------------------------------------------------------
%  CHƯƠNG 3 — PHÂN TÍCH VÀ THIẾT KẾ
% ---------------------------------------------------------------------
\chapter{Phân tích và thiết kế}
\label{ch:thietke}

\begin{muctieu}
Chương này phân tích yêu cầu, thiết kế kiến trúc mạng, thiết kế dữ liệu, bảng
lệnh giao thức và sơ đồ lớp của hệ thống.
\end{muctieu}

\section{Phân tích bài toán và yêu cầu}

Hệ thống có hai loại máy khách: trọng tài gửi kết quả trận đấu và khán giả xem
bảng xếp hạng. Đầu vào là một dòng lệnh \code{SUBMIT} chứa tên hai đội và tỉ
số. Đầu ra là bảng xếp hạng được phát tới mọi khán giả đang xem mỗi khi có trận
mới, kèm các lệnh tra cứu lịch sử đối đầu và xuất bảng ra tệp. Yêu cầu chức
năng gồm sáu lệnh giao thức. Yêu cầu phi chức năng: nhiều trọng tài gửi cùng
lúc không được mất hay trùng trận, một khán giả chậm không được chặn các trọng
tài, và máy chủ khởi động lại phải khôi phục được dữ liệu.

\section{Thiết kế kiến trúc hệ thống}

Kiến trúc client--server được trình bày ở Hình~\ref{fig:cs}. Máy chủ giữ
trạng thái giải trong một đối tượng \code{BangXepHang} đồng bộ. Mỗi kết nối được
phục vụ bởi một luồng riêng từ \code{ExecutorService}. Khán giả gửi
\code{WATCH} để đăng ký vào danh sách phát, sau đó nhận bảng xếp hạng mỗi khi
có trận mới.

\begin{figure}[H]
\centering
\resizebox{\textwidth}{!}{%
\begin{tikzpicture}[node distance=10mm and 16mm]
  \node[objbox] (t1) {Trọng tài 1};
  \node[objbox, below=of t1] (t2) {Trọng tài 2};
  \node[objbox, below=of t2] (tn) {Trọng tài $n$};
  \node[clsbox, right=of t2] (srv) {Máy chủ\\ đa luồng};
  \node[objbox, right=of srv] (k1) {Khán giả 1};
  \node[objbox, above=of k1] (k2) {Khán giả 2};
  \node[objbox, below=of k1] (kn) {Khán giả $m$};
  \draw[flowarr] (t1.east) -- node[above,font=\scriptsize]{SUBMIT} (srv.west);
  \draw[flowarr] (t2.east) -- (srv.west);
  \draw[flowarr] (tn.east) -- (srv.west);
  \draw[flowarr] (srv.east) -- node[above,font=\scriptsize]{STANDINGS} (k2.west);
  \draw[flowarr] (srv.east) -- (k1.west);
  \draw[flowarr] (srv.east) -- (kn.west);
\end{tikzpicture}}
\caption{Sơ đồ client--server: trọng tài gửi kết quả, máy chủ phát bảng xếp hạng.}
\label{fig:cs}
\end{figure}

\section{Bảng lệnh giao thức}

Mỗi lệnh một dòng, phân tách bằng khoảng trắng. Câu trả lời nhiều dòng được
đóng khung bằng một dòng tiêu đề chứa số dòng kế tiếp, ví dụ
\code{STANDINGS <n>}. Bảng~\ref{tab:lenh} liệt kê đầy đủ các lệnh.

\begin{table}[H]
\centering
\caption{Bảng lệnh giao thức.}
\label{tab:lenh}
\begin{tabular*}{\textwidth}{@{\extracolsep{\fill}}>{\bfseries\color{daunavy}}l >{\color{daumaroon}}l l}
\toprule
\rowcolor{gray100}\bfseries\color{daunavy}Lệnh & \bfseries\color{daumaroon}Tham số & \bfseries\color{daunavy}Phản hồi \\
\midrule
WATCH    & -- & Đăng ký xem, gửi bảng hiện tại rồi nhận mỗi lần cập nhật \\
SUBMIT   & độiNhà độiKhách bànNhà bànKhách & \code{OK <trận>}, phát bảng mới \\
STANDINGS& -- & \code{STANDINGS <n>} và $n$ dòng xếp hạng \\
HISTORY  & độiA độiB & \code{HISTORY <n>} và $n$ dòng trận đã gặp \\
EXPORT   & -- & \code{OK <tệp>}, ghi CSV lịch sử và bảng xếp hạng \\
SCHEDULE & các đội (biến thiên, $\geq$ 2) & \code{SCHEDULE <n>} và $n$ dòng lịch vòng tròn \\
QUIT     & -- & \code{BYE}, đóng kết nối \\
\bottomrule
\end{tabular*}
\end{table}

\section{Thiết kế dữ liệu}

Trận đấu là một \code{record} bất biến với xác thực trong hàm tạo compact: tên
đội không rỗng và không chứa khoảng trắng, hai đội khác nhau, tỉ số không âm.
Thống kê của mỗi đội được lưu trong lớp \code{Doi} có các bộ đếm
thắng/hòa/thua và bàn thắng/bàn bại. Bảng xếp hạng là một \code{Map} từ tên đội
sang \code{Doi}, dùng \code{TreeMap} không phân biệt hoa thường để
``HaNoi'' và ``hanoi'' là cùng một đội. Lịch sử trận là một \code{List<Tran>}
theo thứ tự thời gian. Mỗi dòng bảng xếp hạng khi xuất ra là một
\code{record DongXepHang} bất biến, đảm bảo khách nhận được ảnh chụp không bị
sửa khi đang gửi.

\section{Sơ đồ lớp}

Hình~\ref{fig:lop} trình bày sơ đồ lớp chính. Tầng dữ liệu (\code{model}) không
biết gì về mạng; tầng lõi (\code{core}) tính toán xếp hạng; tầng mạng
(\code{server}) chỉ là lớp mỏng gọi vào tầng lõi.

\begin{figure}[H]
\centering
\resizebox{\textwidth}{!}{%
\begin{tikzpicture}[node distance=8mm and 12mm]
  \node[clsbox] (tran) {Tran (record)};
  \node[clsbox, right=of tran] (doi) {Doi};
  \node[clsbox, right=of doi] (dxh) {DongXepHang (record)};
  \node[clsbox, below=of tran] (bxh) {BangXepHang};
  \node[clsbox, right=of bxh] (qt) {QuyTacXepHang};
  \node[clsbox, below=of bxh] (gt) {GiaoThuc};
  \node[clsbox, right=of gt] (lt) {LuuTru};
  \node[clsbox, right=of qt] (mc) {MayChuBXH};
  \node[clsbox, below=of mc] (pk) {PhienKhach};
  \draw[flowarr] (bxh) -- (tran);
  \draw[flowarr] (bxh) -- (doi);
  \draw[flowarr] (doi) -- (dxh);
  \draw[flowarr] (bxh) -- (qt);
  \draw[flowarr] (qt) -- (dxh);
  \draw[flowarr] (mc) -- (bxh);
  \draw[flowarr] (mc) -- (gt);
  \draw[flowarr] (mc) -- (lt);
  \draw[flowarr] (pk) -- (gt);
  \draw[flowarr] (mc) -- (pk);
\end{tikzpicture}}
\caption{Sơ đồ lớp: tầng model, core, protocol, luutru, server.}
\label{fig:lop}
\end{figure}

\section{Thiết kế đồng thời}

Việc nhập một trận gồm ba bước phải nguyên tử: kiểm tra, thêm vào lịch sử, cập
nhật thống kê hai đội. Cả khối được bọc trong \code{synchronized} của
\code{BangXepHang}, nên nhiều trọng tài gửi cùng lúc không thể xen kẽ và làm
hỏng bộ đếm. Việc phát broadcast dùng \code{CopyOnWriteArrayList} cho danh sách
khách xem: thao tác đọc (lặp để gửi) nhiều hơn nhiều so với thêm/xóa, đúng trường
hợp sử dụng của cấu trúc này. Một khoá riêng \code{khoaGui} bảo vệ việc gửi
multi-dòng để hai thông điệp không bị xen dòng vào nhau.

\begin{luuy}
Đặt \code{synchronized} ở mức thao tác nghiệp vụ chứ không phải ở mức từng
trường: nếu không thì hai luồng có thể cùng cập nhật hai đội của cùng một trận
mà không thấy nhau, dẫn đến bảng xếp hạng lệch.
\end{luuy}
